%% kapa: exam=1 round=20 year=2009 period=1 form=A subject=civil_law subject_ko=민법 seq=24 number=24
%% source: Q-Net 감정평가사 기출문제 (cst00309) artlSeq=5024178 file=extracted/1008/5024178/2009년 제20회 감정평가사 1차시험 기출문제/1교시.hwp
%% conversion: kordoc 4.24.51 -> tools/gen1.py (auto); review=proofread-sonnet; flags=none
\begin{question}{24}
공유에 관한 설명으로 옳지 \uline{않은} 것은? (다툼이 있으면 판례에 의함)
\begin{choices}
  \item 소수지분권자가 공유물의 전부 또는 일부를 배타적 독점적으로 사용하는 경우에 다른 공유자는 자신의 지분이 과반수에 미달되더라도 단독으로 공유물 전부의 인도를 청구할 수 있다.
  \item 과반수의 지분권을 가진 공유자가 공유물의 특정부분을 배타적으로 사용ㆍ수익할 것을 정하더라도 이는 공유물의 관리방법으로서 적법하다.
  \item 공유물의 처분ㆍ변경에는 공유자 전원의 동의가 있어야 한다.
  \item 제3자가 공유물을 무단으로 사용한 경우, 공유자 1인은 공유자 전원을 위하여 제3자에게 부당이득 전부의 반환을 청구할 수 있다.
  \item 공유자가 그의 지분을 포기하거나 상속인 없이 사망하면 그 지분은 다른 공유자에게 각 지분의 비율로 귀속한다.
\end{choices}
\end{question}
\answer{4}
